% Source: https://artofproblemsolving.com/community/c3h1222918_glitches_in_math
% Stable identity: aops-topic-1222918-q02
The continued fraction means the limit of $x_1=6$, $x_{k+1}=\dfrac{6}{1+x_k}$. From $x_{k+1}-2=\dfrac{-2(x_k-2)}{1+x_k}$ and $x_{k+1}+3=\dfrac{3(x_k+3)}{1+x_k}$ we get $\dfrac{x_{k+1}-2}{x_{k+1}+3}=-\dfrac{2}{3}\cdot\dfrac{x_k-2}{x_k+3}$, so $\dfrac{x_k-2}{x_k+3}\to0$ and $x_k\to2$.

The equation $x^2+x-6=0$ only says that \emph{if} the limit exists, it is $2$ or $-3$; it does not say which. All $x_k$ are positive, so the limit cannot be $-3$: the value is $x=\boxed{2}$, and $-3$ is an extraneous root.
